% answers_short.tex —— 1977 年·上海卷·文科 · 简短答案速查
% 由 tools/make_short_answers.py 自 answers.tex 的 \gksolhead/\gkanswer 提取，
% 并与 1977/1977_数学真题_上海文_答案要点.md 逐题比对后生成，勿手改（改源 MD 后重跑）。

\begin{gknotice}
\textbf{使用说明}：本节只给最终答案便于快速核对，完整推导见解析版。答案为编写者独立推导，\textbf{非官方标准答案}。
\end{gknotice}

\gkpart{题目}
\begin{itemize}[leftmargin=2.6em]
  \item[\textbf{1.}] (1) $-\dfrac12$；(2) 今年比去年增加 $25\%$；明年养猪 $168$ 头；(3) $x=2$
  \item[\textbf{2.}] $DE=\dfrac{ab}{a+b}$。规范题面按交叉核定采用 $DE\parallel BC$；固定扫描仍保留“$BC$ 的平分线”异文。
  \item[\textbf{3.}] (1) $\dfrac{b-a}{a+b}$；(2) $x<5$
  \item[\textbf{4.}] (1) $\sqrt2+\dfrac{2\sqrt3}{3}$；(2) 见解析版；(3) $BC=4\sqrt6$
  \item[\textbf{5.}] $7251\ \mathrm{mm^3}$（未四舍五入的表达式为 $6000\sqrt3-1000\pi\ \mathrm{mm^3}$）。按通常六角螺帽构形解释。
  \item[\textbf{6.}] 斜率 $k=-\dfrac{\sqrt3}{3}$，倾角为 $150^\circ$。
  \item[\textbf{7.}] $x=4$ 时函数值最小，最小值为 $-11$。
  \item[\textbf{8.}] 取 $340$ 升 $96\%$ 的甲液、$260$ 升 $36\%$ 的乙液。按化学语境将原卷“流酸”理解为硫酸。
\end{itemize}
